
%\documentclass[oneside,final,14pt]{extreport}
\documentclass[12pt, a4paper, openany]{book}
\usepackage[left=2cm,right=2cm,top=1.5cm,bottom=1.5cm,bindingoffset=0cm]{geometry}
%\usepackage[koi8-r]{inputenc}
\usepackage[russianb]{babel}
\usepackage{vmargin}
\setpapersize{A4}
\usepackage[T2A]{fontenc}
\usepackage[utf8x]{inputenc} % more recent versions (at least>=2004-17-10)
%\usepackage[russian]{babel}
%\setmarginsrb{2cm}{1.5cm}{1cm}{1.5cm}{0pt}{0mm}{0pt}{13mm}
\usepackage{indentfirst}
\usepackage{nicefrac} % For comparison
%\usepackage{xfrac} % Works better with other fonts
%\usepackage[unicode, pdftex]{hyperref}
\usepackage{lettrine}
\usepackage[usenames]{color}
%\special{papersize=a5}
\usepackage{colortbl}
\usepackage[pagestyles]{titlesec}
\usepackage{xfrac} % Works better with other fonts
\usepackage[colorlinks=true,linkcolor=black,urlcolor=black,bookmarksopen=true]{hyperref}


% Настройка вертикальных и горизонтальных отступов
\titlespacing{\chapter}{0pt}{5pt}{5pt}
\titlespacing{\section}{\parindent}{4mm}{4mm}
\titlespacing{\subsection}{\parindent}{3mm}{3mm}

\newcommand{\anonsection}[1]{ \section*{#1} \addcontentsline{toc}{section}{\numberline {}#1}} 

\makeatletter %%%%% <---- Starting chapter without a pagebreak
\renewcommand\chapter{\par%
	\thispagestyle{plain}% \global\@topnum\z@
	\@afterindentfalse \secdef\@chapter\@schapter}
\makeatother %%%%% <---- Starting chapter without a pagebreak
\titleformat{\chapter}[display]
{\normalfont\bfseries}{}{0pt}{\Large}

\newpagestyle{mystyle}{
	\sethead[\thepage][][]{}{}{\thepage}	
}

\pagestyle{mystyle}

\sloppy
\begin{document}

	
	\begin{titlepage}
		
		\begin{center}
			%\vfill
			
			%\vfill
			\topskip0pt
			
\begin{flushright}
					{\large\bf АКАДЕМИЯ НАУК СССР\\}
				
				Серия <<История науки и техники>>
\end{flushright}
			
			\vspace*{\fill}
\begin{flushright}
				{\large\bf В. М. ПАСЕЦКИЙ\\}
			\ \\
			\ \\
			{\Huge\bf ПЕРВООТКРЫВАТЕЛИ\\ НОВОЙ\\ ЗЕМЛИ\\}
\end{flushright}

			\vspace*{\fill}
			
			\vfill
			
				\begin{flushright}
			
			ИЗДАТЕЛЬСТВО <<НАУКА>>
		
			Москва 1980
			
		\end{flushright}
		\end{center}
		
	\end{titlepage}
	
	\thispagestyle{empty} % выключаем отображение номера для этой страницы
	
	\newpage
	

\setcounter{secnumdepth}{0}
	
	\section[Введение]{\center ВВЕДЕНИЕ}
	
Открытие и исследование Новой Земли охватывает почти тысячелетие, от первых походов поморов в XI—XII вв. до плаваний в 1910—1911 гг. вокруг северного и южного островов Новой Земли выдающегося полярного исследова-теля Владимира Александровича Русанова, принадлежавшего к первому поколению русских социал-демократов. 

При исследовании исторических источников от русских летописей до архивных документов начала XX в. был собран новый обширный материал об экспедициях в по-лярные страны. Он позволил, хотя бы в самых общих чертах, восстановить картину великих подвигов многих поколений мореплавателей и ученых, которые из века в век стремились к этому исполинскому острову, ревниво охранявшему «льдяными препонами» подступы к таинственному Северо-Восточному проходу. 

Их героическими усилиями медленно приподнималась завеса таинственности, окружавшая Новую Землю. Ее первооткрывателей природа встречала сурово: жестокие морозы, льды и снега, порой почти наглухо заносившие их убогие жилища. Недели, месяцы, годы они питались сухарями и солониной, пили затхлую воду, плавали на 
открытых суденышках среди льдов. 

Документальные материалы, посвященные изучению Новой Земли, весьма обширны. Но они относятся в основном к XIX — началу XX в. Это десятки дневников и описаний путешествий, сотни рапортов и донесений, тысячи научных наблюдений, и в частности метеорологических и магнитных. 

Несравненно более скудны источники, относящиеся к более ранним временам, в особенности к первым столетиям существования Русского государства. Перед историком стоит весьма сложная задача — по отдельным, случайно уцелевшим свидетельствам восстановить картину первых плаваний и первых представлений русских поморов о Новой Земле, западные берега которой они изведали до того, как к ним впервые приблизились англичане, а затем голландцы. Очень важным представляется раскрыть задачи и ход знаменитых путешествий Виллема Баренца, со своими спутниками зимовавшего на северо-восточной оконечности Новой Земли, в Ледяной гавани. Решение этой задачи облегчается тем, что отдел истории Нидерландов Государственного музея в Амстердаме предоставил в распоряжение автора гравюры и карты, относящиеся к плаваниям Баренца. 

Существенную часть этого исторического повествования составляет рассказ о первых русских государственных экспедициях на Новую Землю и плавании помора Саввы Лошкина, впервые обошедших все берега Новой Земли — от Карских Ворот до мыса Желания. 

Особое внимание уделено плаваниям Г. В . Поспелова, Ф. П. Литке, П. К . Пахтусова, А. К . Цивольки, усилиями которых были достоверно картированы почти все западные и восточные берега Новой Земли. Благодаря их исследованиям на смену полуфантастическим представлениям об этом острове, якобы соединявшемся с Америкой, пришли точные знания и объективные научные представления не только о берегах Новой Земли, но и о ее природе, в том числе первые инструментальные геофизические наблюдения как на берегах Новой Земли, так и в ее водах. Именно их труды проложили путь к первым попыткам освоения великой северной морской трассы и сделали возможным в дальнейшем правильное судоходство не только в Баренцевом, но и в Карском море. 

Подвижническая деятельность Г. В. Поспелова, Ф. П. Литке, П. К. Пахтусова, А. К. Цивольки показана в основном на базе новых документальных источников, с использованием научного и эпистолярного наследства В. М. Головнина, Г. А. Сарычева, И. Ф. Крузенштерна, Ф. П. Литке, М. Ф. Рейнеке и др. 

Особенно важно было всесторонне раскрыть деятельность первой академической экспедиции на Новую Землю, которую возглавлял великий естествоиспытатель К. М. Бэр. Это тем более необходимо, что в научной и научно-популярной литературе бытует неверное утверждение, что Бэр назвал Карское море ледником и тем самым якобы задержал освоение западного участка Северного морского пути. В действительности, анализируя первые метеорологические наблюдения, он прозорливо предсказал существование Новоземельского ледяного массива, который является серьезным препятствием даже для современных судов. Бэр к тому же наметил план широкого изучения Севера в естественнонаучном отношении, который обнаружен нами в Центральном государственном историческом архиве в Ленинграде. 

Не менее существенным представляется раскрытие исследований России на Новой Земле во время Первого международного полярного года и научных предприятий Академии наук, в том числе плавания академика А. Ф. Миддендорфа, геологических изысканий академика Ф. Н. Чернышева, наблюдений академика Б. Б. Голицына во время солнечного затмения 1896 г. Заключительная часть книги посвящена путешествиям на Новую Землю Владимира Александровича Русанова, который первым из ученых обошел сначала ее северный, а затем и южный остров, повторив тем самым подвиг помора Саввы Лошкина. Русанов пять лет посвятил исследованию Новой Земли, составил карты ее северных берегов, создал цикл работ по геологии и географии этого арктического архипелага. Его подвиг — одна из самых ярких, самых удивительных страниц истории изучения Арктики. С новоземельскими исследованиями связано его решение отправиться в плавание по Северному морскому пути на деревянном «Геркулесе», который бесследно исчез во льдах вместе со своими отважными обитателями. Судьба этой экспедиции и по сей день привлекает внимание многих советских людей. 

И наконец, на последних страницах этой книги будет рассказано об открытиях и наблюдениях на Новой Земле экспедиции Г. Я. Седова, которая на пути к Северному полюсу зазимовала у берегов северного острова Новой Земли, и ее участники пересекли покрывавший его исполинский ледяной купол. 

Это небольшое историческое повествование в значительной степени основано на новых архивных материалах, собранных в течение четверти века и составляющих лишь часть обширнейшего комплекса нередко уникальных документов об открытиях и исследованиях в Арктике с древнейших времен до начала XX в. 

В этом поиске добрую помощь автору постоянно оказывали сотрудники Ленинградского отделения архива Академии наук СССР, Центрального государственного архива Военно-Морского Флота, Центрального государственного исторического архива в Ленинграде, Центрального государственного архива древних актов, Центрального военно-исторического архива, Центрального государственного исторического архива ЭССР в Тарту, Центрального государственного архива Октябрьской революции, Государственных архивов Новгородской, Орловской, Архангельской областей, Музея Арктики и Антарктики, Центрального музея Военно-Морского Флота, библиотеки Ленинградского отделения Института истории естествознания и техники АН СССР, библиотеки Ленинградского отделения Института истории СССР АН СССР, библиотеки Главной геофизической обсерватории им. А. И. Воейкова, Государственной публичной библиотеки им. М. Е. Салтыкова-Щедрина и Государственной библиотеки им. В. И. Ленина. 

	\section[Первые плавания к Новой Земле]{\center ПЕРВЫЕ ПЛАВАНИЯ К НОВОЙ ЗЕМЛЕ }
	
	\newpage
	\tableofcontents
	
	\thispagestyle{empty} % 
	
	\newpage
	
	\setcounter{secnumdepth}{0}
	
	\phantomsection
	
		\section*{Описание}
	
	{\bf Название:} Первооткрыватели Новой Земли
	
{\bf Автор:} Пасейкий Василий Михайлович
	
{\bf Издательство:} Москва: «Наука» 1980
	
		{\bf Ответственный редактор:} \textit{академик А. П. ОКЛАДНИКОВ}

	
		{\bf Аннотация:} В книге рассказывается об истории открытия од-ного из самых обширных архипелагов Арктики. От  первого похода помора Саввы Лошкина до плавания  вокруг Новой Земли знаменитого полярного путешественника Владимира Александровича Русанова, таков путь новоземельских изысканий. 
		
		Большое внимание уделено исследованиям В. Баренца, Г.В. Поспелова, Ф.П. Литке, П.К. Пахтусова, А.К. Цивольки, доставивших науке первые сведения о фауне и флоре островов, их геологическом строении, климате, промысловых богатствах.
		\thispagestyle{empty} % выключаем отображение номера для этой страницы

	
\end{document}


